\scanpage{9}
\uline{Thm 2} Pour $\forall\ \pi$ discret, il existe un $\pi$ fibré universel,
i.e. un espace classifiant (défini à isom unique près dans la cat. homotopique \begin{tikzpicture}[x=.25cm,y=.25cm,baseline=-.4ex]\fill (0,0) circle (.08);\draw (.3,-.02)..controls (.55,.06) and (.62,.04)..(.7,-.02)..controls (.8,-.11) and (.95,-.11)..(1.05,-.1);\end{tikzpicture})

Soit $X$ un espace sur lequel $\pi$ opère à \unsure{g.}
et prenons

\diagram{\begin{tikzcd}[ampersand replacement=\&,column sep=large]
\&E_\pi\times_\pi X\overset{\deleted{\unsure{E|df}}}{=}\underbrace{(E_\pi\times X)/\pi}\arrow[dl]\arrow[dr]\&\supsetX\\
B_\pi=E_\pi/\pi\&\&X/\underset{\deleted{E_\pi}}{\pi}
\end{tikzcd}}

\margin{\begin{tikzpicture}[scale=.45]\draw[gray] (0,0)--(.8,.35)--(.6,.65) (.8,.35)--(.4,.35) (.4,-.1)--(.42,-.4);\end{tikzpicture}}
\diagram{\begin{tikzcd}[ampersand replacement=\&,column sep=small]
\pi_2(B_\pi)\arrow[r]\arrow[d,equal]\&\pi_1(X)\arrow[r,dash]\&\pi_1(X,\pi)\arrow[r]\&\pi\arrow[r]\&\pi_0(\unsure{V|X})\arrow[d,equal]\\
0\&\&\&\&1
\end{tikzcd}}

\mspage{8}
